\documentclass[11pt]{article}
\usepackage{amsmath,amssymb,amsthm}
\usepackage[margin=1in]{geometry}
\usepackage[T1]{fontenc}
\title{Is $\|S\|_{L^\infty}\le C\|\omega\|_{L^\infty}$? A deep check of \texttt{PeriodicStrainLinfBound}}
\author{Benjamin Stanley Frohman}
\date{2 November 2026}
\newtheorem{theorem}{Theorem}\newtheorem{proposition}{Proposition}
\theoremstyle{definition}\newtheorem{definition}{Definition}
\begin{document}\maketitle
\begin{center}Copyright \copyright\ 2026 Benjamin Stanley Frohman. Released under CC BY 4.0.\\
Lean record \texttt{PeriodicNSContinuation/StrainLinf.lean} (Apache-2.0).\\
This document does not claim a solution of the Navier--Stokes Millennium Problem.\end{center}

\begin{abstract}
\texttt{PeriodicStrainLinfBound} asks for one constant $C$ with $\|S(t,x)\|\le C\|\omega(t)\|_{L^\infty}$ for every smooth periodic Navier--Stokes solution. We show it is false as stated, already at $t=0$: an explicit family of divergence-free trigonometric polynomials on $\mathbb{T}^3$ has $\|\omega_K\|_\infty\le 1.403$ and $\|S_K\|_\infty\ge\tfrac2\pi\log K$. The arithmetic divergence is proved in Lean; the Fourier identity and the Gibbs bound are classical and are checked numerically. We record the conditional results that do give the bound, and what remains of it for Navier--Stokes solutions.
\end{abstract}

\section{Literature}
Biot--Savart gives $\nabla u=\mathcal{R}\omega$ with $\mathcal{R}$ a matrix of Calder\'on--Zygmund operators (compositions of Riesz transforms). These are bounded on $L^p$, $1<p<\infty$, and from $L^\infty$ to BMO, but not on $L^\infty$. For bounded vorticity the velocity is only log-Lipschitz (Yudovich; see Majda--Bertozzi, \emph{Vorticity and Incompressible Flow}, Ch.~8). The standard sharp example is the Bahouri--Chemin steady state $\omega=\operatorname{sgn}x\operatorname{sgn}y$ on $\mathbb{T}^2$, whose velocity gradient diverges logarithmically at the origin (Bahouri--Chemin 1994; used by Denisov for double-exponential growth of $\nabla\omega$ in 2D Euler, and by Elgindi--Jeong in 3D symmetric settings). The quantitative substitute is the log inequality of Beale--Kato--Majda and Kozono--Taniuchi, $\|\nabla u\|_\infty\le C(1+\|\omega\|_{\mathrm{BMO}}(1+\log^+\|u\|_{H^s}))$.

\section{An explicit periodic counterexample family}
Let $f_K(x)=\tfrac4\pi\sum_{k\ \mathrm{odd}\le K}\sin(kx)/k$, the Dirichlet partial sum of $\operatorname{sgn}(\sin x)$. On $\mathbb{T}^3$ put
\[
\omega_K=f_K(x)f_K(y)\,e_z,\qquad \Delta\psi_K=f_Kf_K,\qquad u_K=(-\partial_y\psi_K,\ \partial_x\psi_K,\ 0).
\]
Then $u_K$ is a divergence-free trigonometric polynomial with $\nabla\times u_K=\omega_K$. With $\omega_K=\sum A_{kl}\sin kx\sin ly$, $A_{kl}=16/(\pi^2kl)$, we have $\psi_K=-\sum A_{kl}\sin kx\sin ly/(k^2+l^2)$ and
\[
S_{11}(0)=-\partial_{xy}\psi_K(0)=\frac{16}{\pi^2}\sum_{k,l\ \mathrm{odd}\le K}\frac1{k^2+l^2},\qquad S_{12}(0)=0,
\]
so $\|S_K(0)\|_{op}=S_{11}(0)$.
\begin{theorem}[\texttt{bcSum\_unbounded}, proved in Lean]
$\sum_{i,j<n}\frac{1}{(2i+1)^2+(2j+1)^2}\to\infty$. More precisely it is at least $\tfrac18\sum_{i<n}\tfrac1{i+1}$.
\end{theorem}
The Gibbs bound $\sup|f_K|\le \tfrac2\pi\operatorname{Si}(\pi)\approx1.1790$, uniform in $K$, is classical; so $\|\omega_K\|_\infty\lesssim1.39$ while $\|S_K\|_\infty\ge S_{11}(0)\sim\tfrac2\pi\log K$. The ratio $\|S\|_\infty/\|\omega\|_\infty$ is invariant under $u\mapsto\lambda u(\mu\,\cdot)$, so no rescaling removes the divergence.

\paragraph{Numerics (complete; exact sums at the origin, grid sup of $\|S\|_{op}$ on an $8K\times8K$ grid up to $K=255$).}
\begin{center}\small\begin{tabular}{rrrrr}$K$&$\|\omega_K\|_\infty$&$S_{11}(0)$&$\sup_x\|S\|_{op}$&$S_{11}(0)-\tfrac2\pi\log K$\\\hline
7 & 1.4024 & 1.6585 & 1.6585 & 0.4196 \\
15 & 1.3931 & 2.0977 & 2.0977 & 0.3737 \\
31 & 1.3908 & 2.5385 & 2.5385 & 0.3523 \\
63 & 1.3902 & 2.9796 & 2.9796 & 0.3420 \\
127 & 1.3900 & 3.4209 & 3.4209 & 0.3370 \\
255 & 1.3900 & 3.8621 & 3.8621 & 0.3345 \\
511 & 1.3900 & 4.3034 & -- & 0.3332 \\
1023 & 1.3900 & 4.7447 & -- & 0.3326 \\
\end{tabular}\end{center}
\begin{center}\small\begin{tabular}{rrrrr}$K$&$t$&$M$&$\|S\|_\infty$&ratio\\\hline
31 & 0.00 & 1.3908 & 2.5385 & 1.8252 \\
31 & 0.20 & 1.0757 & 2.2023 & 2.0473 \\
31 & 0.40 & 1.0199 & 1.9882 & 1.9494 \\
31 & 0.60 & 1.0075 & 1.8391 & 1.8253 \\
31 & 0.80 & 1.0042 & 1.7235 & 1.7163 \\
31 & 1.00 & 1.0023 & 1.6293 & 1.6255 \\
127 & 0.00 & 1.3899 & 3.4209 & 2.4612 \\
127 & 0.20 & 1.0000 & 2.2628 & 2.2628 \\
127 & 0.40 & 1.0000 & 2.0156 & 2.0156 \\
127 & 0.60 & 1.0000 & 1.8563 & 1.8563 \\
127 & 0.80 & 1.0000 & 1.7350 & 1.7350 \\
127 & 1.00 & 1.0000 & 1.6371 & 1.6371 \\
\end{tabular}\end{center}

The first table shows $\|\omega_K\|_\infty\to1.3900=(\tfrac2\pi\operatorname{Si}\pi)^2$, the sup of $\|S\|$ attained at the origin, and $S_{11}(0)-\tfrac2\pi\log K\to\approx0.332$: growth exactly $\tfrac2\pi\log K$. The second table evolves the $K=31$ and $K=127$ data by 2D Navier--Stokes ($=$ $z$-independent 3D NS), $\nu=0.005$, to $t=1$: the ratio is largest near $t=0$ and is then quickly smoothed (to $\approx1.63$ by $t=1$).

\section{Conditional results that do give the bound}
\begin{itemize}
\item \textbf{Dini continuity} (classical CZ theory): if $\omega(t)$ has modulus of continuity $\varpi$ with $\int_0^1\varpi(r)\,dr/r<\infty$, then $\|\nabla u\|_\infty\le C\bigl(\|\omega\|_\infty+\int_0^1\varpi(r)\frac{dr}r+\|\omega\|_{L^2}\bigr)$. This includes H\"older and log-H\"older moduli $\varpi(r)\lesssim|\log r|^{-a}$, $a>1$. The constant depends on the modulus, not on $\|\omega\|_\infty$ alone; this is a different statement.
\item \textbf{Log inequality} (BKM, Kozono--Taniuchi): bounded by $\|\omega\|_\infty$ or $\|\omega\|_{\mathrm{BMO}}$ times $\log$ of a higher norm. This is the right replacement and is what Section~2 of the cubic-source note now uses.
\item \textbf{Direction regularity} (Constantin--Fefferman): Lipschitz (or $\tfrac12$-H\"older, Beir\~ao da Veiga--Berselli) continuity of $\xi=\omega/|\omega|$ on the high-vorticity set controls the stretching $\xi\cdot S\xi$ through a local, cancelling kernel. It does not give $\|S\|_\infty\le C\|\omega\|_\infty$; the family above is two-dimensional, with constant direction $\xi=e_z$, and still violates the bound.
\end{itemize}

\section{Could it hold for Navier--Stokes solutions?}
\begin{enumerate}
\item \textbf{Uniform constant: false.} \texttt{PeriodicStrainLinfBound} quantifies over all smooth periodic data with one $C$, at all $t\ge0$ including $t=0$. Each $u_K$ is smooth periodic divergence-free data with a smooth local solution (classical local well-posedness; it is in fact global, being 2D). At $t=0$ the bound fails as $K\to\infty$. In Lean this is not yet a formal refutation: the repo has no local-existence theorem to produce \texttt{IsPeriodicMaximalSolution} for $u_K$, so the Prop remains an unassumed named hypothesis, now recorded as false modulo classical local existence.
\item \textbf{Constant depending on the solution, up to $T^*$: open, and not enough.} Any smooth solution satisfies $\|S\|_\infty\le C_T\|\omega\|_\infty$ on each $[0,T]$, $T<T^*$, trivially. Whether $C$ can be taken uniform up to a finite $T^*$ is open. Even if it could, it gives only $M'\le CM^2$, a lifespan bound, not regularity (cubic-source note, \S3).
\item \textbf{Positive-time smoothing.} Viscosity regularises the data instantly, so a bound with $C$ depending on $\nu$, $t>0$ and the energy might hold away from $t=0$ (the numerics are consistent with this). Near a putative singularity this is exactly what is unknown.
\end{enumerate}

\section{Record}
\textbf{Proved (Lean):} \texttt{bcSum\_unbounded}. \textbf{Classical, checked numerically:} Fourier identity for $S_{11}(0)$, Gibbs bound, $\tfrac2\pi\log K$ growth. \textbf{Conclusion:} \texttt{PeriodicStrainLinfBound} as stated is false (modulo local existence, not formalized). \textbf{Open:} solution-dependent constant up to $T^*$; not sufficient for regularity in any case.
\end{document}
